\documentclass[11pt]{letter}

\usepackage[a4paper, margin=1in]{geometry}
\usepackage[T1]{fontenc}
\usepackage{lmodern}
\usepackage{hyperref}
\hypersetup{colorlinks=true, urlcolor=black, linkcolor=black}

\setlength{\parindent}{0pt}
\setlength{\parskip}{10pt}

\begin{document}

\begin{center}
{\Large \textbf{Ivan Evdokimov, Ph.D.}}\\[4pt]
{\small Manchester, UK \; $|$ \; ivevdm@gmail.com \; $|$ \; \href{https://linkedin.com/in/ivan-evdokimov}{linkedin.com/in/ivan-evdokimov} \; $|$ \; \href{https://github.com/ivan020}{github.com/ivan020}}
\end{center}

\vspace{10pt}

Dear Hiring Manager,

I'm writing to apply for the Research Data Scientist role at dunnhumby. I hold a PhD in Computational Finance from the University of Essex and have spent the last several years building and deploying machine learning systems in production, most recently on classification and risk-assessment models for large-scale datasets. I'm excited about this role's focus on insight automation and product assortment, and about the chance to apply generative AI to product development in a market-leading Customer Data Science business.

What I'd bring to this role is a consistent track record of taking an open-ended problem, researching the right approach, and landing on a solution that's both correct and scalable in practice. At UK Data Service, I trained and fine-tuned PyTorch and scikit-learn models for automated metadata classification and disclosure risk assessment, reaching 95\%+ F1 accuracy on sensitive social science datasets, then built the ETL pipelines needed to run them against PostgreSQL-hosted data at 100K+ record scale. When the existing disclosure-control algorithms couldn't keep up with growing data volumes, I researched alternative approaches, migrated the core logic from R to C++ using bitmask techniques, and delivered an 11x performance improvement that enabled near real-time processing on datasets exceeding 500K records. That pattern --- diagnosing why an approach doesn't scale, researching alternatives, and engineering a solution that holds up under real production load --- is exactly the kind of science-to-module work this role is built around.

That same instinct extends to how I work with new tools and problem spaces more broadly: rather than defaulting to a familiar technique, I look for the approach that actually fits the data and constraints in front of me, then build it out properly rather than settling for something that only works in a notebook. I've presented research at conferences and taught data science and machine learning to 100+ students with a 4.8/5 feedback average, so communicating technical reasoning clearly to both technical and non-technical audiences is something I do regularly, not just something I can do.

I'm particularly drawn to dunnhumby's focus on turning data science into reusable, client-facing solutions, and to the opportunity to explore generative AI applications for product attribute generation and product development. I'd welcome the chance to discuss how my background in applied ML and large-scale data pipelines could contribute to your team.

Thank you for your consideration.

Kind regards,\\
Ivan Evdokimov

\end{document}
